\section{DBの開始}

\subsection{実施にかかる手続き（労使合意、承認認可）}


\begin{InyouJobunbox}
  \begin{itemize}
    \item 法第3条
    \item 則第2条〜第4条、第11条
    \item 解釈第8-2、8-7、8-1
  \end{itemize}
\end{InyouJobunbox}


\begin{fillblank}[自作 2026/10/4]

  \noindent \textbf{◯DB法}
  \begin{lawstructure}
    \begin{lawarticle}{第三条}{確定給付企業年金の実施}
      厚生年金適用事業所の事業主は、確定給付企業年金を実施しようとするときは、確定給付企業年金を実施しようとする厚生年金適用事業所に使用される（　　A　　）で組織する労働組合があるときは当該労働組合、当該（　　A　　）で組織する労働組合がないときは当該（　　A　　）を代表する者の同意を得て、確定給付企業年金に係る規約（以下「規約」という。）を作成し、次の各号のいずれかに掲げる手続を執らなければならない。
      \lawitem{一}当該規約について厚生労働大臣の承認を受けること。
      \lawitem{二}企業年金基金（以下「基金」という。）の設立について厚生労働大臣の認可を受けること。

      \lawparagraph{2}
      確定給付企業年金は、一の厚生年金適用事業所について一に限り実施することができる。ただし、政令で定める場合においては、この限りでない。
      \lawparagraph{3}
      二以上の厚生年金適用事業所について確定給付企業年金を実施しようとする場合においては、第一項の同意は、（　　B　　）について得なければならない。
    \end{lawarticle}
  \end{lawstructure}

\end{fillblank}

\begin{answerbox}
 \begin{enumerate}[A:]
    \item 厚生年金保険の被保険者の過半数
    \item 各厚生年金適用事業所
  \end{enumerate}
\end{answerbox}


\begin{fillblank}[自作 2026/10/4]

  \noindent \textbf{◯DB則}
  \begin{lawstructure}
    \begin{lawarticle}{第三条}{過半数代表者}
      法第三条第一項、法第六条第二項（法第七条第二項において準用する場合を含む。）及び法第七十八条第一項の規定による手続を厚生年金保険の被保険者（法第二条第三項に規定する厚生年金保険の被保険者をいう。以下同じ。）の過半数を代表する者（以下この条において「過半数代表者」という。）の同意を得て行う場合にあっては、当該過半数代表者は、次の各号のいずれにも該当する者とする。
      \lawitem{一}労働基準法（昭和二十二年法律第四十九号）第四十一条第二号に規定する（　　A　　）の地位にある者でないこと。
      \lawitem{二}過半数代表者を選出することを明らかにして実施される（　　B　　）等の方法による手続により選出された者であって、（　　C　　）に基づき選出されたものでないこと。
      
      \lawparagraph{2}
      前項第（　　※出題の都合で略　　）号に該当する者がいない厚生年金適用事業所にあっては、過半数代表者は同項（　　D　　）に該当する者とする。
      \lawparagraph{3}
      確定給付企業年金を実施しようとする又は実施する厚生年金適用事業所の事業主は、当該事業主に使用される者が過半数代表者であること若しくは過半数代表者になろうとしたこと又は過半数代表者として正当な行為をしたことを理由として（　　E　　）をしないようにしなければならない。
      \lawparagraph{4〜5} （略）
    \end{lawarticle}
  \end{lawstructure}

\end{fillblank}

\begin{answerbox}
 \begin{enumerate}[A:]
    \item 監督又は管理
    \item 投票、挙手
    \item 事業主の意向
    \item 第二号
    \item 不利益な取り扱い
  \end{enumerate}
\end{answerbox}



\begin{truefalse}[年金1 \ 2019 \ 問題1(2)\textcircled{6}]
\textbf{下線部の内容の正誤を判定し、誤っている場合は正しい内容に改めなさい。}
\\

確定給付企業年金法施行規則第3 条において、規約の作成、規約の変更および実施事業所
の増減等に関する同意手続きを、厚生年金保険の被保険者の過半数を代表する者の同意を得
て行う場合には、当該者は次の A、B のいずれにも該当する者とすることとされている。
\begin{enumerate}[A]
  \item 労働基準法第41 条第2 号に規定する監督または管理の地位にある者でないこと
  \item 過半数代表者を選出することを明らかにして実施される投票、挙手等の方法による手続
きにより選出された者\underline{であること}
\end{enumerate}
\end{truefalse}

\begin{answerbox}
 \textbf{誤} \\
  〜により選出された者\underline{であって、事業主の意向により選出された者でないこと}
\end{answerbox}

\begin{truefalse}[自作 \ 2026/10/4]
\textbf{以下の文章の正誤を判定し、誤っている場合はその理由を記述しなさい。}
\\

厚生年金適用事業所の事業主は、確定給付企業年金を実施しようとするとき、
確定給付企業年金を実施しようとする厚生年金適用事業所に使用される厚生年金保険の被保険者の過半数で
組織する労働組合があるときであっても、
当該厚生年金保険の被保険者の過半数を代表する者の同意さえ得られれば、確定給付企業年金を実施できる。
\end{truefalse}

\begin{answerbox}
 \textbf{誤} \\
  「過半数組合の同意」か「過半数代表者の同意」かの選択は、過半数組合の有無によって自動で決まるものであり、
  事業主が任意に選択できるものではない。
\end{answerbox}

\begin{explanationbox}
  出典：DB法第三条
\end{explanationbox}

\newpage

\subsection{１事業所１DBの原則}



\begin{InyouJobunbox}
  \begin{itemize}
    \item 法第3条第2項
    \item 令第1条
    \item 則第1条
  \end{itemize}
\end{InyouJobunbox}


\begin{fillblank}[自作 2026/10/4]

  \noindent \textbf{◯DB令}
  \begin{lawstructure}
    \begin{lawarticle}{第一条}{複数の確定給付企業年金を実施できる場合}
      確定給付企業年金法（以下「法」という。）第三条第二項ただし書の政令で定める場合は、一の厚生年金適用事業所（法第二条第二項に規定する厚生年金適用事業所をいう。以下同じ。）について二の確定給付企業年金を実施する場合であって当該二の確定給付企業年金のうちいずれか一方の確定給付企業年金を実施する厚生年金適用事業所（以下「実施事業所」という。）の事業主（第五十三条並びに附則第三条及び第八条を除き、以下「事業主」という。）の全部が同時に（　　A　　）とならないときその他厚生労働省令で定める場合とする。
    \end{lawarticle}
  \end{lawstructure}

  \noindent \textbf{◯DB則}
  \begin{lawstructure}
    \begin{lawarticle}{第一条}{複数の確定給付企業年金を実施できる場合}

      確定給付企業年金法施行令（平成十三年政令第四百二十四号。以下「令」という。）第一条の厚生労働省令で定める場合は、次のとおりとする。

      \lawitem{一}
      一の厚生年金適用事業所（確定給付企業年金法（平成十三年法律第五十号。以下「法」という。）第二条第二項に規定する厚生年金適用事業所をいう。以下同じ。）について二以上の確定給付企業年金を実施する場合であって、それぞれの確定給付企業年金の加入者（以下「加入者」という。）について適用される（　　B　　）、（　　C　　）その他これらに準ずるもの（以下「（　　B　　）等」という。）が異なる場合
      \lawitem{二}
      法人である確定給付企業年金を実施する事業主（第三条第一項第二号、第三項及び第五項、第十九条の二第二号イ、第百二十条、附則第六条第一項第一号、附則第七条第一項並びに附則第十二条第一項第一号を除き、以下「事業主」という。）が他の法人である事業主と（　　D　　）した場合であって、当該（　　D　　）の日から起算して（　　E　　）を経過していない場合
      \lawitem{三}
      給付の額の算定方法が第二十五条第四号に掲げる方法である確定給付企業年金（以下「（　　F　　）」という。）と（　　F　　）でない確定給付企業年金とをそれぞれ実施する場合
    \end{lawarticle}
  \end{lawstructure}

\end{fillblank}


\begin{answerbox}
 \begin{enumerate}[A:]
    \item 他方の確定給付企業年金の事業主の全部
    \item 労働協約
    \item 就業規則
    \item 合併
    \item 原則として一年
    \item リスク分担型企業年金
  \end{enumerate}
\end{answerbox}


\begin{truefalse}[年金1 \ 2019 \ 問題1(2)\textcircled{3}]
\textbf{下線部の内容の正誤を判定し、誤っている場合は正しい内容に改めなさい。}
\\

例外的に2 つ以上の確定給付企業年金を実施することが認められているケースとして、会
社合併した場合であって、\underline{当該合併にかかる規約変更の日から起算して2 年}を経過していな
い場合が該当する。
\end{truefalse}

\begin{answerbox}
 \textbf{誤} \\
  〜合併した場合であって、\underline{当該合併の日から起算して原則として1年}を経過していない場合〜
\end{answerbox}



\begin{truefalse}[年金1 \ H29 \ 問題2(1)\textcircled{3}]
\textbf{以下の文章の正誤を判定し、誤っている場合その理由を記述しなさい。}
\\

リスク分担型企業年金とリスク分担型企業年金でない確定給付企業年金の経理をそれぞれで
行うことを規約に定めることで、１つの規約型確定給付企業年金において、リスク分担型企業
年金とリスク分担型企業年金でない確定給付企業年金を併存することができる。
\end{truefalse}

\begin{answerbox}
 \textbf{誤} \\
  経理をそれぞれ行うとともに、資産をそれぞれ区分して運用することを規約に定める必要がある。
\end{answerbox}


\begin{explanationbox}
  参照：承認認可基準別紙第3-2(4) 

  \textgt{
  リスク分担型企業年金とリスク分担型企業年金でない確定給付企業年金の経理をそれぞれで行うとともに、資産をそれぞれに区分して運用することを規約に定める措置。なお、基金型の場合においては、2-4(6)の事項を規約に定める措置もあわせて講じること。
  }
\end{explanationbox}



\begin{shortanswer}[年金1 \ 2023 \ 問題2(1)ア]
１つの厚生年金適用事業所が複数の確定給付企業年金を実施することができる場合として、
確定給付企業年金法施行令第１条に定める場合（注）\underline{以外}の、確定給付企業年金法施行規則
第１条に定める場合を３つ簡潔に入力しなさい。 （250 字以内）\\
（注）複数の厚生年金適用事業所が共同で実施する確定給付企業年金に加入する一方で、企業
独自に実施する確定給付企業年金に加入する等、企業年金を実施する事業主の範囲が異
なる場合。
\end{shortanswer}

\begin{answerbox}
 \begin{itemize}
  \item それぞれの確定給付企業年金の加入者について適用する労働協約、就業規則その他これらに準ずるものが異なる場合
  \item 法人である確定給付企業年金を実施する事業主が他の法人である事業主と合併した場合であって、当該合併した日から起算して原則として一年を経過していない場合
  \item リスク分担型企業年金とリスク分担型企業年金でない確定給付企業年金とをそれぞれ実施する場合
 \end{itemize}
\end{answerbox}

\begin{explanationbox}
  類題：年金2 2022 問題2(3)
\end{explanationbox}